\documentclass[a4paper, 10pt, ]{article}

\input{../../KUT000/preamble_v1.0.tex}

% -----------------------------------------------------------------------------

\def\oznacenieCelku{Kolekcia učebných textov}

% -----------------------------------------------------------------------------


\def\KUTporadoveCislo{022}

\def\oznacenieVerzie{v0.96}
% \def\oznacenieVerzie{\phantom{v1.0}}

\def\mesiacRok{marec 2026}

\def\authorslabel{AS, MT}





% -----------------------------------------------------------------------------

\begin{document}

% -----------------------------------------------------------------------------
% Uvodny nadpis

\noindent
\parbox[t][18mm][c]{0.3\textwidth}{%
\raisebox{-0.9\height}{%
\phantom{.}
\includegraphics[height=18mm]{../../KUT000/URKFEIlogo_v0.1.pdf}%
}%
}%
\parbox[t][18mm][c]{0.7\textwidth}{%
\raggedleft

\sffamily
\fontsize{16pt}{18pt}
\fontseries{sbc}
\selectfont

\noindent
\textcolor[rgb]{0.75, 0.75, 0.75}{\textls[25]{\oznacenieCelku}}
}%

\noindent
\parbox[t][16mm][b]{0.5\textwidth}{%
\raggedright

\color{Gray}
\sffamily

\fontsize{12pt}{12pt}
\selectfont
\mesiacRok

\fontsize{6pt}{10pt}
\selectfont
github.com/OkoliePracovnehoBodu/KUT

\fontsize{8pt}{10pt}
\selectfont
\authorslabel




}%
\parbox[t][16mm][b]{0.5\textwidth}{%
\raggedleft

\sffamily

\fontsize{6pt}{6pt}
\selectfont

\textcolor[rgb]{0.68, 0.68, 0.68}{\oznacenieVerzie}


\fontsize{14pt}{14pt}
\selectfont

\bfseries


\includegraphics[height=12pt]{../../KUT000/KUT_logo_v0.1.pdf}%
{%
\textls[-50]{\KUTporadoveCislo}
}%
}%

% -----------------------------------------------------------------------------




\vspace{6mm}

% ---------------------------------------------
\sffamily
\bfseries
\fontsize{18pt}{21pt}
\selectfont

\begin{flushleft}
    TS: Príklad spracovania datasetu pre prevodovú charakteristiku
\end{flushleft}

\bigskip

% -----------------------------------------------------------------------------
\normalsize
\normalfont
% -----------------------------------------------------------------------------

\lstset{style=mystyle}










\noindent
\lettrine[lines=1, nindent=1pt, loversize=0.0]{C}{ieľom} 
textu je uvedenie príkladu spracovania dát, z ktorých je možné získať prevodovú charakteristiku daného dynamického systému. Tento text priamo nadväzuje na \href{run:../../KUT021/TeX/KUT021.pdf}{\textsf{KUT021}}, kde je prezentovaný práve súbor takýchto dát.

 






\section{Získané dáta}


Príklad získaných dát pre systém TS pri prevádzkovej podmienke $6$~[V] je znázornený na nasledujúcom obrázku \ref{fj_01_data_ts_pot6}.

\begin{figure}[!b]

\vbox{%
    \makebox[\textwidth][c]{%
    \includegraphics{exp3_meranie.png}
    }

    \figcaption{
        Získané dáta pre vstupný signál a výstupné signály systému TS pri manuálne nastavenej prevádzkovej podmienke ventilátora $6$~[V].
    }
    \label{fj_01_data_ts_pot6}
}%vbox

\end{figure}

Dátový súbor obsahuje informácie o ustálených stavoch systému TS pre zvolené hodnoty vstupného signálu $u_1$ (napätie ovládajúce výhrevné teleso) a prevádzkovej podmienky $u_2$ (napätie ovládajúce ventilátor). V experimentoch boli realizované merania pre vstupný signál $u_1$ v rozsahu od $0$~[V] do $10$~[V] s~krokom $1$~[V] a~pre prevádzkovú podmienku $u_2$ v hodnotách uvedených v tabuľke \ref{tab:ustalene_hodnoty_podmienok_TS}. Ustálené hodnoty vstupného signálu $u_1$ sú uvedené v tabuľke \ref{tab:ustalene_hodnoty_vstupneho_signalu_TS}.




\begin{figure}[!t]

\vbox{%
    \makebox[\textwidth][c]{%
    \includegraphics{exp1_meranie.png}
    }

    \figcaption{
        Získané dáta systému TS pri prevádzkovej podmienke ventilátora $2$~[V], pri ktorých dochádza k prechodu systému do havarijného režimu.
    }
    \label{fj_01_data_ts_emergency}
}%vbox

\end{figure}



\begin{table}[t!]

    \vbox{%
    \tabcaption{Ustálené hodnoty vstupného signálu (výhrevné teleso) [V]}
    \label{tab:ustalene_hodnoty_vstupneho_signalu_TS}

    \vspace{-10pt}

    \lstyle

    \begin{tabular*}{\textwidth}{@{\extracolsep{\fill}} ccccccccccc}
    \toprule
    0 & 1 & 2 & 3 & 4 & 5 & 6 & 7 & 8 & 9 & 10 \\
    \bottomrule
    \end{tabular*}
    }%vbox

\end{table}



\begin{table}[t!]

    \vbox{%
    \tabcaption{Ustálené hodnoty prevádzkovej podmienky (ventilátor) [V]}
    \label{tab:ustalene_hodnoty_podmienok_TS}

    \vspace{-10pt}

    \lstyle

    \begin{tabular*}{\textwidth}{@{\extracolsep{\fill}} ccccc}
    \toprule
    2 & 4 & 6 & 8 & 10 \\
    \bottomrule
    \end{tabular*}
    }%vbox

\end{table}





\subsection{Výskyt neštandardných meraní}

Dáta na obrázku \ref{fj_01_data_ts_pot6} poskytujú informácie o ustálených stavoch systému pri prevádzkovej podmienke $6$~[V] a pre rozsah vstupného signálu od $0$~[V] do $10$~[V].

Systém TS je vybavený ochranou proti prekročeniu maximálnej teploty. Ak senzory indikujú dosiahnutie limitnej hodnoty, ventilátory sa automaticky prepnú na maximálny výkon a vypne sa výhrevné teleso, zariadenie prechádza do havarijného režimu.
Príklad merania, pri ktorom zariadenie prechádza do havarijného režimu, je znázornený na obrázku \ref{fj_01_data_ts_emergency}. V tomto experimente bola prevádzková podmienka ventilátora nastavená na hodnotu $2$~[V].


V experimente je možné pozorovať, že pri vstupnom napätí špirály približne od $7$~[V] do $10$~[V] dochádza k aktivácii tejto ochrany. V tejto oblasti preto systém nedosahuje ustálený stav a tieto merania sa pri ďalšej analýze neuvažujú ako ustálené stavy.











\section{Spracovanie dát pre prevádzkovú podmienku 6 [V]}

\subsection{Vyznačenie ustálených úsekov}

Zo získaných dát na obrázku \ref{fj_01_data_ts_pot6} je potrebné vyňať úseky, počas ktorých sú výstupné a~vstupné signály ustálené. V tomto prípade nech za ustálený stav sa považuje posledných $30$ sekúnd zvoleného časového intervalu pre ustálenie, pričom posledných $0.1$ sekundy merania sa neuvažuje. Takto vyňaté úseky sú znázornené na obrázku \ref{fj_02_data_ts_pot6}.


\begin{figure}[!t]

\vbox{%
    \makebox[\textwidth][c]{%
    \includegraphics{exp3_meranie_steadyPoints.png}
    }

    \vspace{-6pt}

    \caption{
        %
    }
    \label{fj_02_data_ts_pot6}
}%vbox

\vspace{-10pt}

\end{figure}









\subsection{Nameraná prevodová charakteristika}

Vyňaté úseky je možné zlúčiť do jedného dátového súboru. Obsahuje len hodnoty vstupného signálu a výstupných signálov, pre ktoré je možné konštatovať, že zodpovedajú ustálenému stavu. Takéto dáta charakterizujú systém v ustálených stavoch.

Grafické znázornenie uvedených dát tak, že na x-ovej osi sú hodnoty vstupného signálu v ustálenom stave a na y-ovej osi sú prislúchajúce hodnoty výstupných signálov zo snímača 1 a snímača 2 v ustálenom stave, je grafickým vyjadrením nameranej prevodovej charakteristiky systému. Nameraná prevodová charakteristika je znázornená na obrázku \ref{fj_03_data_ts_panel_1}.

Pre lepšiu ilustráciu, že uvedený dataset neobsahuje len niekoľko bodov, ale že každý ustálený stav je reprezentovaný viacerými hodnotami, je na obrázku \ref{fj_03_data_ts_panel_2} zobrazený detail pre zvolenú hodnotu vstupného signálu.

% \begin{center}
\begin{figure}[!b]

\vspace{-12pt}

\vbox{%
    \makebox[\textwidth][c]{%
    \includegraphics{exp3_prevodova.png}
    }

    \vspace{-6pt}

    \caption{
        Nameraná prevodová charakteristika systému TS pre prevádzkovú podmienku $6$ [V].
    }
    \label{fj_03_data_ts_panel_1}
}%vbox

\vspace{-10pt}

\end{figure}
% \end{center}




% \begin{center}
\begin{figure}[!t]

\vbox{%
    \makebox[\textwidth][c]{%
    \includegraphics{exp3_prevodova_u1.png}
    }

    \vspace{-6pt}

    \figcaption{
        %
    }
    \label{fj_03_data_ts_panel_2}

    \vspace{-6pt}

}%vbox

\end{figure}
% \end{center}





\begin{figure}[!b]

\vbox{%
    \makebox[\textwidth][c]{%
    \includegraphics{exp3_prevodova_fit_sat.png}
    }

    \vspace{-6pt}

    \figcaption{
        Aproximácia nameranej prevodovej charakteristiky systému TS lineárnym modelom so saturovanou časťou.
    }
    \label{fj_04_data_ts_panel_1}
}%vbox

\end{figure}





\subsection{Aproximácia prevodovej charakteristiky}
\label{sec:Aproximácia prevodovej charakteristiky}


Nameraná prevodová charakteristika znázornená na obrázku \ref{fj_03_data_ts_panel_1} zachytáva vzťah medzi vstupným signálom a výstupnými signálmi systému v ustálených stavoch. Výstup systému je v tomto prípade reprezentovaný dvoma meraniami zo snímača~1 a snímača~2, ktoré označme ako $y_1$ a $y_2$.

Grafické znázornenie ukazuje, že pre hodnoty vstupného signálu približne od $7$~[V] dochádza k saturovaniu výstupných veličín. Z tohto dôvodu je účelné rozdeliť prevodovú charakteristiku na dve časti.

Vstupnú veličinu označme ako $u$. Závislosť výstupných veličín od vstupnej veličiny potom je
\begin{equation}
y_1 = f_1(u), \qquad y_2 = f_2(u)
\end{equation}

Na úseku $0 \leq u \leq 7$~[V] je možné na základe grafického zobrazenia nameraných dát aproximovať závislosť lineárnou funkciou. Parametre lineárnej aproximácie boli určené metódou najmenších štvorcov pomocou funkcie \texttt{polyfit} z knižnice \texttt{NumPy}. Výsledkom sú modely
\begin{align}
y_1 &= 0,709533\,u + 3,951781 \\
y_2 &= 0,559831\,u + 4,135048
\end{align}


Pre hodnoty vstupného signálu $7 < u \leq 10$~[V] sú výstupné veličiny saturované. Priemerné hodnoty výstupov v tejto oblasti sú
\begin{align}
y_1 &= 8,500301 \\
y_2 &= 7,815181
\end{align}

Takto získaná aproximácia prevodovej charakteristiky je znázornená na obrázku \ref{fj_04_data_ts_panel_1}.



\begin{figure}[!t]

\vbox{%
    \makebox[\textwidth][c]{%
    \includegraphics{exp3_prevodova_fit_sat_poly2.png}
    }

    \vspace{-6pt}

    \figcaption{
        Aproximácia prevodovej charakteristiky systému TS polynómom druhého stupňa so saturovanou časťou.
    }
    \label{fj_04_data_ts_panel_2}

    \vspace{-10pt}

}%vbox

\end{figure}


\begin{figure}[!b]

    \vspace{-10pt}

    \vbox{%
    \makebox[\textwidth][c]{%
    \includegraphics{allExp_prevodova.png}
    }

    \vspace{-6pt}

    \caption{
        Prevodové charakteristiky systému TS pre rôzne prevádzkové podmienky ventilátora.
    }
    \label{fj_05_ts_panel_1}

    \vspace{-10pt}

    }%vbox

\end{figure}

 


Okrem lineárnej aproximácie bola zostrojená aj aproximácia polynómom druhého stupňa. Parametre modelu boli opäť určené pomocou funkcie \texttt{polyfit}. Pre úsek $0 \leq u \leq 7$~[V] majú modely tvar
\begin{align}
y_1 &= -0,046524\,u^2 + 1,035201\,u + 3,626113 \\
y_2 &= -0,025066\,u^2 + 0,735293\,u + 3,959586
\end{align}

Aj v tomto prípade je pre rozsah $7 < u \leq 10$~[V] výstup modelu reprezentovaný konštantnou hodnotou zodpovedajúcou saturovanej oblasti merania.

Ak by sme chceli model vyjadriť ako celok, formálne je možné písať pre lineárnu aproximáciu
\begin{align}
\hat y_1 &=
\begin{cases}
0,709533\,u + 3,951781 & \text{pre } 0 \leq u \leq 7 \\
8,500301 & \text{pre } 7 < u \leq 10
\end{cases}
\\
\hat y_2 &=
\begin{cases}
0,559831\,u + 4,135048 & \text{pre } 0 \leq u \leq 7 \\
7,815181 & \text{pre } 7 < u \leq 10
\end{cases}
\end{align}
Analogicky je možné zapísať aj model založený na polynóme druhého stupňa
\begin{align}
\hat y_1 &=
\begin{cases}
-0,046524\,u^2 + 1,035201\,u + 3,626113 & \text{pre } 0 \leq u \leq 7 \\
8,500301 & \text{pre } 7 < u \leq 10
\end{cases}
\\
\hat y_2 &=
\begin{cases}
-0,025066\,u^2 + 0,735293\,u + 3,959586 & \text{pre } 0 \leq u \leq 7 \\
7,815181 & \text{pre } 7 < u \leq 10
\end{cases}
\end{align}












\section{Spracovanie dát pre všetky prevádzkové podmienky}

Obdobne ako na obr.~\ref{fj_03_data_ts_panel_1} je možné vizualizovať všetky namerané prevodové charakteristiky pre všetky prevádzkové podmienky (p.p.), pozri obrázok \ref{fj_05_ts_panel_1}. 

Na tieto dáta, a teda na skúmaný systém ako celok, je možné nazerať ako na dáta zodpovedajúce systému s dvomi vstupmi a dvomi výstupmi. Prvým vstupom je vstupný signál $u$ (napätie špirály), druhým vstupom je prevádzková podmienka $u_2$ (napätie ventilátora). Výstupy systému sú reprezentované hodnotami zo snímača~1~a~snímača~2, ktoré označme ako $y_1$ a~$y_2$.

Dáta reprezentujúce ustálené stavy je potom možné vizualizovať ako funkciu dvoch premenných. Takéto grafické znázornenie je uvedené na obrázku \ref{fj_06_ts_panel_1}.













\subsection{Aproximácia s využitím viacfaktorového modelu}

Na základe uvedených dát je možné zostrojiť model systému, ktorý závisí od dvoch vstupných veličín. Model bol určený metódou najmenších štvorcov s využitím dát zo všetkých experimentov.

Model lineárneho typu má tvar
\begin{align}
y_1 &= -0,142576\,u_2 + 0,770612\,u + 4,815384 \\
y_2 &= -0,051529\,u_2 + 0,595590\,u + 4,375411
\end{align}
Grafické porovnanie takéhoto modelu s nameranými hodnotami je znázornené na obrázku \ref{fj_07_ts_panel_1}.



Okrem lineárneho modelu je možné zostrojiť aj model založený na polynóme druhého stupňa. Pre výstupy systému majú modely tvar
\begin{align}
\begin{aligned}
y_1 &=
0,025350\,u_2^2
-0,061827\,u_1^2
\\&-0,039657\,u_2 u_1
-0,312671\,u_2
+1,447598\,u_1
+4,275240
\end{aligned}
\\
\begin{aligned}
y_2 &=
-0,001441\,u_2^2
-0,032007\,u_1^2
\\&-0,023818\,u_2 u_1
+0,047715\,u_2
+0,964436\,u_1
+3,615530
\end{aligned}
\end{align}
Grafické znázornenie takéhoto modelu je uvedené na obrázku \ref{fj_07_ts_panel_2}.


\section{Záverečné poznámky}

Skripty pre spracovanie dát a generovanie obrázkov, ktoré sú uvedené v texte vyššie, sa nachádzajú v podadresári \texttt{PY}, konkrétne v súboroch \texttt{dataJob.ipynb} a \texttt{figJob.ipynb}.








% \begin{figure}[!b]
\begin{center}

% \vspace{-10mm}

\vbox{%
    \makebox[\textwidth][c]{%
    \includegraphics{allExp_3D_static_characteristic.png}
    }

    \vspace{-8mm}

    \figcaption{
        Vizualizácia ustálených stavov systému TS v trojrozmernom priestore. Čiernou farbou sú znázornené hodnoty pre vstupný signál $0 \leq u \leq 7$~[V], červenou farbou hodnoty v saturovanej oblasti $8 \leq u \leq 10$~[V]. Červené krížiky označujú merania, pri ktorých systém prešiel do havarijného režimu.
    }
    \label{fj_06_ts_panel_1}

}%vbox

\end{center}
% \end{figure}




% \begin{figure}[p]
\begin{center}

\vspace{-10mm}

\vbox{%
    \makebox[\textwidth][c]{%
    \includegraphics{allExp_3D_static_characteristic_with_plane.png}
    }

    \vspace{-8mm}

    \figcaption{
        Porovnanie lineárneho viacfaktorového modelu s nameranými dátami. Fialová plocha predstavuje model určený metódou najmenších štvorcov.
    }
    \label{fj_07_ts_panel_1}

    }%vbox

\end{center}
% \end{figure}

\vfill

% \begin{figure}[p]
\begin{center}

    \vspace{-10mm}

    \vbox{%
    \makebox[\textwidth][c]{%
    \includegraphics{allExp_3D_static_characteristic_with_quad_surface.png}
    }

    \vspace{-8mm}

    \figcaption{
        Porovnanie kvadratického viacfaktorového modelu s nameranými dátami. Fialová plocha predstavuje model určený metódou najmenších štvorcov.
    }
    \label{fj_07_ts_panel_2}

    \vspace{20pt}

    }%vbox

\end{center}
% \end{figure}












% -----------------------------------------------------------------------------

\end{document}

% -----------------------------------------------------------------------------